\chapter{Diseases of the Endocrine System}
\label{ch:endocrine}
The endocrine system consists of glands that produce hormones---chemical messengers that regulate metabolism, growth and development, reproduction, and homeostasis. Major endocrine organs include the thyroid, parathyroids, adrenal glands, pituitary gland, and endocrine pancreas (discussed in Chapter~\ref{ch:pancreas}). Endocrine disorders include conditions of hormone excess (hyperfunction), hormone deficiency (hypofunction), and autoimmune and neoplastic diseases.

%-------------------------------------------------------------
\section{Adrenal Diseases}
%-------------------------------------------------------------

%-------------------------------------------------------------

%-------------------------------------------------------------

%-------------------------------------------------------------

Adrenal diseases involve secretory dysfunction or neoplastic growth affecting the adrenal cortex or adrenal medulla. Cortical pathology disrupts steroidogenesis of glucocorticoids, mineralocorticoids, and adrenal androgen precursors, whereas medullary lesions alter catecholamine synthesis. Hormonal excess or deficiency syndromes exert profound systemic effects on blood pressure homeostasis, fluid balance, carbohydrate metabolism, and stress response.

%-------------------------------------------------------------

%-------------------------------------------------------------

%-------------------------------------------------------------

%-------------------------------------------------------------

\label{sec:Adrenal Diseases}
%-------------------------------------------------------------%-------------------------------------------------------------

\subsection{Addison's Disease}
\label{sec:Addison's Disease}
Addison's disease (primary adrenal insufficiency) is destruction or dysfunction of the adrenal cortex, leading to deficiency of cortisol, aldosterone, and adrenal androgens. The most common cause in developed countries is autoimmune adrenalitis (often associated with other autoimmune diseases in APS type 2), while tuberculosis remains the most common cause worldwide; other causes include infections (CMV in HIV/AIDS, fungal), adrenal hemorrhage (Waterhouse-Friderichsen syndrome---meningococcal sepsis), metastatic disease, and drugs (ketoconazole, etomidate). Clinical features include fatigue, weight loss, anorexia, nausea, vomiting, abdominal pain, hyperpigmentation (from elevated ACTH and MSH stimulating melanocytes, a hallmark of primary adrenal insufficiency), salt craving, hypotension, and orthostatic dizziness. Diagnosis shows low morning cortisol, elevated ACTH (primary), and failure to suppress cortisol with high-dose dexamethasone. Management is glucocorticoid replacement (hydrocortisone 15--25 mg/day in divided doses) and mineralocorticoid replacement (fludrocortisone 0.05--0.2 mg/day). Prognosis is good with appropriate replacement therapy and education about stress dosing.

\subsection{Adrenal Incidentaloma}
\label{sec:Adrenal Incidentaloma}
Adrenal incidentalomas are adrenal masses discovered incidentally on imaging performed for unrelated reasons, found in approximately 4\% of CT scans. Most are benign, non-functioning cortical adenomas, but they may represent functioning tumors (cortisol-producing adenoma causing subclinical Cushing's syndrome, aldosterone-producing adenoma, pheochromocytoma) or rarely, adrenocortical carcinoma or metastases. The condition results from incidental detection of adrenal masses. Evaluation requires hormonal workup (cortisol after 1 mg dexamethasone suppression test, plasma metanephrines, aldosterone/renin ratio) and imaging characteristics (size, density, enhancement pattern). Management includes surgery for suspicious masses ($>$4 cm, HU $>$10, irregular borders, rapid growth), functioning tumors, and suspected malignancy, with surveillance for indeterminate masses. Prognosis is excellent for benign non-functioning adenomas.

\subsection{Cushing's Syndrome}
\label{sec:Cushing's Syndrome}
Cushing's syndrome results from chronic exposure to excess cortisol, regardless of the cause, with exogenous glucocorticoid use (iatrogenic) being the most common cause overall. Endogenous causes are divided into ACTH-dependent (70--80\%): Cushing's disease (pituitary adenoma secreting ACTH, 70\%) and ectopic ACTH syndrome (small cell lung cancer, bronchial carcinoid); and ACTH-independent (20--30\%): adrenal adenoma, adrenal carcinoma, and bilateral adrenal hyperplasia. Clinical features include central (truncal) obesity, moon facies, dorsocervical fat pad (``buffalo hump''), supraclavicular fat pads, thin extremities, purple striae ($>$1 cm wide), easy bruising, proximal muscle weakness, hypertension, hyperglycemia/diabetes, osteoporosis, menstrual irregularity, hirsutism, and psychiatric disturbances. Diagnosis involves confirmation of hypercortisolism (24-hour urinary free cortisol, late-night salivary cortisol, or 1 mg overnight dexamethasone suppression test) followed by determination of the cause (plasma ACTH, high-dose dexamethasone suppression test, CRH stimulation test, MRI pituitary, CT adrenals, inferior petrosal sinus sampling). Management depends on the cause: transsphenoidal surgery for Cushing's disease, surgical resection of adrenal tumors, chemotherapy for adrenal carcinoma, and ketoconazole or metyrapone for medical management. Prognosis is favorable with successful treatment.

\subsection{Pheochromocytoma}
\label{sec:Pheochromocytoma}
Pheochromocytoma is a catecholamine-secreting tumor arising from chromaffin cells of the adrenal medulla, accounting for 0.1--0.6\% of hypertensive patients, with paragangliomas being extra-adrenal catecholamine-secreting tumors. Approximately 10\% are bilateral, 10\% are extra-adrenal, and 10\% are malignant, with up to 40\% being hereditary (MEN2, von Hippel-Lindau, neurofibromatosis type 1, SDH mutations). Clinical features include the classic triad of episodic headache, sweating, and palpitations with paroxysmal or sustained hypertension. Diagnosis is by plasma free metanephrines or 24-hour urinary metanephrines and catecholamines (high sensitivity), with imaging including CT or MRI of the abdomen and MIBG scan. Management is alpha-adrenergic blockade (phenoxybenzamine or doxazosin) for at least 10--14 days preoperatively, followed by beta-blockade (only after adequate alpha-blockade), volume expansion, and laparoscopic adrenalectomy. Prognosis is excellent after resection for benign disease.

%-------------------------------------------------------------
\section{Parathyroid Diseases}
%-------------------------------------------------------------

%-------------------------------------------------------------

%-------------------------------------------------------------

%-------------------------------------------------------------

%-------------------------------------------------------------

%-------------------------------------------------------------

%-------------------------------------------------------------

%-------------------------------------------------------------

\label{sec:Parathyroid Diseases}
%-------------------------------------------------------------%-------------------------------------------------------------

\subsection{Hyperparathyroidism}
\label{sec:Hyperparathyroidism}
Hyperparathyroidism is characterized by excessive secretion of parathyroid hormone (PTH). Primary hyperparathyroidism (most common) results from a parathyroid adenoma (85\%), hyperplasia (10--15\%), or rarely carcinoma ($<$1\%), causing hypercalcemia which may be asymptomatic or symptomatic: ``bones, stones, abdominal groans, and psychic moans'' (bone pain/pathologic fractures, kidney stones, abdominal pain/pancreatitis/constipation, depression/cognitive dysfunction). Secondary hyperparathyroidism occurs in CKD due to phosphate retention and decreased vitamin D activation. Tertiary hyperparathyroidism occurs after prolonged secondary hyperparathyroidism with autonomous PTH secretion. Diagnosis of primary disease shows elevated calcium with elevated or inappropriately normal PTH. Management is parathyroidectomy for symptomatic patients or those meeting criteria (calcium $>$1 mg/dL above normal, reduced bone density, kidney stones, age $<$50, or renal impairment). Prognosis is excellent with surgical cure.

\subsection{Hypoparathyroidism}
\label{sec:Hypoparathyroidism}
Hypoparathyroidism is characterized by deficient PTH production, leading to hypocalcemia and hyperphosphatemia. The most common cause is accidental damage to or removal of the parathyroid glands during thyroid surgery (iatrogenic hypoparathyroidism), with other causes including autoimmune hypoparathyroidism (in isolation or as part of APS-1, AIRE gene mutation) and genetic causes such as DiGeorge syndrome (22q11.2 deletion). Clinical features of hypocalcemia include perioral and digital paresthesias, carpopedal spasm, tetany (Chvostek's sign---facial muscle twitch with facial nerve tapping; Trousseau's sign---carpopedal spasm with blood pressure cuff inflation), seizures, and cardiac arrhythmias (prolonged QT). Diagnosis is by low PTH with hypocalcemia and hyperphosphatemia. Management is oral calcium supplementation (calcium carbonate or citrate) and active vitamin D analogs (calcitriol). Prognosis is good with appropriate supplementation.

\subsection{Multiple Endocrine Neoplasia}
\label{sec:Multiple Endocrine Neoplasia}
Multiple endocrine neoplasia syndromes are autosomal dominant disorders characterized by tumors in two or more endocrine glands. MEN1 (Wermer syndrome) involves the parathyroids (primary hyperparathyroidism, most common), pancreas (gastrinoma, insulinoma, non-functional tumors), and pituitary (prolactinoma most common), caused by menin gene mutations. MEN2A (Sipple syndrome) involves medullary thyroid carcinoma (nearly 100\% penetrance), pheochromocytoma, and primary hyperparathyroidism, caused by RET proto-oncogene mutations. MEN2B includes medullary thyroid carcinoma, pheochromocytoma, mucosal neuromas, marfanoid habitus, and intestinal ganglioneuromatosis. The condition results from germline mutations in tumor suppressor genes or proto-oncogenes. Clinical features depend on the glands involved. Genetic testing guides prophylactic thyroidectomy in RET mutation carriers (within the first year of life for MEN2B, by age 5 for MEN2A). Prognosis depends on early detection and intervention.

%-------------------------------------------------------------
\section{Pituitary Diseases}
%-------------------------------------------------------------

%-------------------------------------------------------------

%-------------------------------------------------------------

%-------------------------------------------------------------

Pituitary diseases encompass functional hypersecretion, hormone deficiency syndromes, and mass lesions affecting the anterior and posterior pituitary gland. As the master regulator of the endocrine axis, pituitary pathology alters thyroidal, adrenal, gonadal, and somatic growth signaling. Compressive mass lesions in the sella turcica additionally cause optic chiasm compression, visual field deficits, and cranial nerve palsies.

%-------------------------------------------------------------

%-------------------------------------------------------------

%-------------------------------------------------------------

%-------------------------------------------------------------

\label{sec:Pituitary Diseases}
%-------------------------------------------------------------%-------------------------------------------------------------

\subsection{Acromegaly}
\label{sec:Acromegaly}
Acromegaly is a chronic disorder caused by excessive growth hormone (GH) secretion, almost always from a pituitary somatotroph adenoma (95\%), with GH stimulating hepatic production of insulin-like growth factor 1 (IGF-1) which mediates most of the somatic effects, and onset being insidious with a mean delay of 7--10 years between symptom onset and diagnosis. Clinical features include coarsening of facial features, prognathism (jaw protrusion), frontal bossing, enlarged hands and feet, soft tissue swelling, macroglossia, excessive sweating, deepening of the voice, arthropathy, carpal tunnel syndrome, hypertension, diabetes mellitus, sleep apnea, cardiomyopathy, and increased risk of colorectal polyps and cancer. Diagnosis shows elevated age- and sex-matched IGF-1 level and failure to suppress GH below 1 ng/mL after 75-g oral glucose tolerance test, with MRI showing a pituitary macroadenoma in most cases. Management includes transsphenoidal surgery (first-line, curative in 80--90\% of microadenomas), somatostatin analogs (octreotide LAR, lanreotide, pasireotide) for persistent disease, pegvisomant (GH receptor antagonist), and dopamine agonists (cabergoline) for mild cases. Prognosis is improved with early treatment, which reduces cardiovascular morbidity and mortality.

\subsection{Diabetes Insipidus}
\label{sec:Diabetes Insipidus}
Diabetes insipidus is a condition of impaired water homeostasis caused by deficient production (central DI) or action (nephrogenic DI) of antidiuretic hormone (ADH/vasopressin). Central DI results from decreased ADH secretion due to pituitary/hypothalamic disease (tumors, surgery, trauma, infiltrative disease, idiopathic), while nephrogenic DI results from kidney resistance to ADH (lithium use is the most common acquired cause; genetic forms include X-linked aquaporin-2 mutations). Clinical features include polyuria (output $>$3 L/day, often 5--20 L/day), polydipsia, nocturia, and dilute urine (specific gravity $<$1.005, osmolality $<$300 mOsm/kg), with severe dehydration and hypernatremia possibly occurring if water intake is restricted. Diagnosis involves the water deprivation test, with desmopressin administration differentiating central DI (improved urine concentration) from nephrogenic DI (no response). Management of central DI is desmopressin (DDAVP) replacement, while management of nephrogenic DI includes thiazide diuretics, low-salt diet, and avoidance of nephrotoxic agents. Prognosis is good with appropriate management.

\subsection{Hyperprolactinemia}
\label{sec:Hyperprolactinemia}
Hyperprolactinemia is the most common pituitary hormone excess disorder, characterized by elevated serum prolactin levels. Causes include prolactin-secreting pituitary adenomas (prolactinomas, most common), medications (antipsychotics, metoclopramide, domperidone, verapamil), hypothalamic/pituitary stalk compression, primary hypothyroidism (increased TRH stimulates prolactin), chronic renal failure, and chest wall trauma. Clinical features in women include galactorrhea, oligomenorrhea or amenorrhea, infertility, and decreased libido. In men: decreased libido, erectile dysfunction, infertility, gynecomastia, and galactorrhea (rare). Mass effects from macroprolactinomas (>1 cm) include headache, visual field defects (bitemporal hemianopia), and hypopituitarism. Diagnosis is based on confirmed elevated serum prolactin (macroprolactin excluded), MRI pituitary with contrast to identify adenoma, and thyroid function tests to exclude hypothyroidism. Management includes dopamine agonists (cabergoline preferred over bromocriptine due to better efficacy and tolerability) as first-line therapy for symptomatic hyperprolactinemia and prolactinomas, and transsphenoidal surgery for dopamine agonist-resistant or apoplectic macroprolactinomas. Prognosis is excellent; most patients achieve normal prolactin levels with dopamine agonist therapy.

\subsection{Hypopituitarism}
\label{sec:Hypopituitarism}
Hypopituitarism is deficiency in one or more pituitary hormones, with causes including pituitary adenomas (most common), craniopharyngiomas, meningiomas, pituitary surgery/radiation, Sheehan syndrome (postpartum pituitary necrosis from hemorrhage), traumatic brain injury, infiltrative diseases (sarcoidosis, histiocytosis), infections, and autoimmune hypophysitis, with hormones lost in a characteristic order: GH $>$ LH/FSH $>$ TSH $>$ ACTH $>$ prolactin. Clinical manifestations depend on which hormones are deficient: GH deficiency (decreased muscle mass, increased fat), gonadotropin deficiency (amenorrhea, erectile dysfunction, decreased libido), ACTH deficiency (secondary adrenal insufficiency---without hyperpigmentation and with preserved aldosterone), TSH deficiency (central hypothyroidism), and ADH deficiency (central diabetes insipidus). Diagnosis requires dynamic endocrine testing (insulin tolerance test, ACTH stimulation test, GnRH stimulation test) and MRI of the pituitary. Management is hormone replacement for each deficient axis, with hydrocortisone for ACTH deficiency (must be replaced before thyroid hormone to avoid adrenal crisis). Prognosis is good with appropriate replacement.

\subsection{SIADH}
\label{sec:SIADH}
Syndrome of inappropriate antidiuretic hormone secretion (SIADH) is the most common cause of euvolemic hyponatremia, characterized by excessive ADH secretion despite normal or low plasma osmolality. Causes include malignancies (small cell lung cancer most common, producing ectopic ADH), pulmonary disorders (pneumonia, tuberculosis, aspergillosis), CNS disorders (meningitis, encephalitis, stroke, traumatic brain injury), and medications (SSRIs, TCAs, carbamazepine, vincristine, cyclophosphamide). Clinical features depend on the severity and rate of development of hyponatremia: mild (lethargy, confusion, headache), moderate (nausea, vomiting, gait instability), severe (seizures, coma, respiratory arrest). Diagnosis requires hyponatremia (<135 mEq/L), low plasma osmolality (<275 mOsm/kg), inappropriately concentrated urine (urine osmolality >100 mOsm/kg), elevated urine sodium (>40 mEq/L), and exclusion of hypovolemia, hypervolemia, adrenal insufficiency, and hypothyroidism. Management includes fluid restriction (first-line), pharmacotherapy (demeclocycline, vaptans/tolvaptan for refractory cases), and for severe symptomatic hyponatremia, cautious 3\% hypertonic saline with frequent sodium monitoring to avoid osmotic demyelination syndrome. Prognosis relates to the underlying cause; SIADH itself resolves when the causative condition is treated.

%-------------------------------------------------------------
\section{Thyroid Diseases}
%-------------------------------------------------------------

%-------------------------------------------------------------

%-------------------------------------------------------------

%-------------------------------------------------------------

Thyroid diseases represent a broad spectrum of structural and functional disorders of the thyroid gland, ranging from benign nodules and diffuse goiters to thyroiditis and overt neoplasia. Dysregulation of thyroid hormone production alters basal metabolic rate, thermogenesis, cardiovascular dynamics, and neuromuscular irritability. Evaluation combines sensitive serum immunoassay panels, high-resolution thyroid ultrasonography, fine-needle aspiration biopsy, and radionuclide imaging.

%-------------------------------------------------------------

%-------------------------------------------------------------

%-------------------------------------------------------------

%-------------------------------------------------------------

\label{sec:Thyroid Diseases}
%-------------------------------------------------------------%-------------------------------------------------------------

\subsection{Addison Disease}
\label{sec:Addison Disease}
Addison disease (primary adrenal insufficiency) is chronic adrenocortical destruction leading to deficiency of cortisol, aldosterone, and adrenal androgen production. The condition results from autoimmune adrenalitis in Western nations (80\% of cases, often associated with APS-1 or APS-2) and tuberculosis worldwide. Clinical features include fatigue, weight loss, hyperpigmentation of skin folds and palmar creases (due to elevated ACTH and MSH cleavage from POMC), postural hypotension, hyponatremia, hyperkalemia, and cravings for salt. Adrenal crisis is a life-threatening acute presentation with severe hypotension, vomiting, and confusion. Diagnosis is confirmed by low morning serum cortisol with elevated ACTH level, and failure to stimulate cortisol in response to the cosyntropin (synthetic ACTH) test. Management requires lifelong glucocorticoid replacement (hydrocortisone) and mineralocorticoid replacement (fludrocortisone). Prognosis is excellent with adequate replacement therapy.

\subsection{Cushing Syndrome}
\label{sec:Cushing Syndrome}
Cushing syndrome is a clinical state resulting from prolonged exposure to excess systemic glucocorticoids. Exogenous administration of glucocorticoids is the most common cause. Endogenous Cushing syndrome results from Cushing disease (ACTH-secreting pituitary adenoma, 70\% of endogenous cases), ectopic ACTH secretion (e.g., small cell lung carcinoma), or primary adrenal adenoma/carcinoma. Clinical features include central obesity, moon facies, buffalo hump, abdominal purple striae, proximal muscle weakness, skin fragility, hypertension, hyperglycemia, and osteopenia. Diagnosis involves initial screening with 24-hour urinary free cortisol, late-night salivary cortisol, or 1-mg overnight dexamethasone suppression test, followed by plasma ACTH measurement to differentiate ACTH-dependent from ACTH-independent causes. Management is targeted: surgical resection of pituitary or adrenal tumors, discontinuation of exogenous steroids, or medical therapy (ketoconazole, metyrapone). Prognosis is good following successful surgical cure.

\subsection{Graves Disease}
\label{sec:Graves Disease}
Graves disease is an autoimmune disorder and the most common cause of hyperthyroidism (60--80\% of cases), occurring predominantly in women aged 20--50. The condition results from thyroid-stimulating immunoglobulin (TSI/TSHR autoantibodies) binding to and activating the TSH receptor, stimulating unregulated thyroid hormone synthesis and diffuse thyroid hyperplasia. Clinical features include heat intolerance, weight loss despite increased appetite, tachycardia, palpitations, fine tremor, diffuse non-tender goiter with bruit, and specific extrathyroidal manifestations: Graves ophthalmopathy (exophthalmos, proptosis, lid lag due to orbital fibroblast TSHR expression) and pretibial myxedema. Diagnosis is confirmed by suppressed TSH, elevated free T4 and T3, positive TSI antibodies, and diffuse increased radioactive iodine uptake (RAIU). Management includes antithyroid drugs (methimazole, propylthiouracil), radioactive iodine ablation, or thyroidectomy, along with beta-blockers (propranolol) for symptom control. Prognosis is generally favorable.

\subsection{Hashimoto Thyroiditis}
\label{sec:Hashimoto Thyroiditis}
Hashimoto thyroiditis (chronic lymphocytic thyroiditis) is an autoimmune disease and the most common cause of hypothyroidism in iodine-sufficient regions, affecting women 10 times more frequently than men. The condition results from cell-mediated cytotoxicity (CD8+ T cells) and anti-thyroid peroxidase (anti-TPO) and anti-thyroglobulin autoantibodies, leading to lymphocytic infiltration, germinal center formation, Hurthle cell metaplasia, and progressive destruction of thyroid follicles. Clinical features include cold intolerance, weight gain, fatigue, constipation, dry skin, bradycardia, muscle weakness, and a non-tender firm goiter. Patients may experience transient hyperthyroidism (hashitoxicosis) early in the disease due to follicular destruction. Diagnosis is confirmed by elevated serum TSH, low free T4, and high titers of anti-TPO antibodies. Management is lifelong hormone replacement with levothyroxine, adjusted to maintain TSH within normal limits. Prognosis is excellent with proper dosing.

\subsection{Hyperthyroidism}
\label{sec:Hyperthyroidism}
Hyperthyroidism (thyrotoxicosis) is a state of excessive thyroid hormone production, affecting approximately 1--2\% of the population. The most common cause is Graves' disease (an autoimmune condition with thyroid-stimulating immunoglobulins [TSI] that activate the TSH receptor, causing diffuse goiter, ophthalmopathy, and dermopathy), with other causes including toxic multinodular goiter, toxic adenoma (Plummer's disease), subacute thyroiditis (painful, viral), silent thyroiditis, postpartum thyroiditis, and exogenous thyroid hormone ingestion. Clinical features include weight loss, heat intolerance, palpitations, anxiety, tremor, diarrhea, menstrual irregularities, warm moist skin, lid retraction, and diffuse goiter. Diagnosis shows suppressed TSH with elevated free T4 and/or T3. Management options include antithyroid drugs (methimazole---preferred, or propylthiouracil [PTU]---used in first trimester of pregnancy and thyroid storm), radioactive iodine (I-131) ablation, and thyroidectomy, with beta-blockers (propranolol) providing symptomatic relief. Prognosis is good with appropriate treatment.

\subsection{Hypothyroidism}
\label{sec:Hypothyroidism}
Hypothyroidism is a state of deficient thyroid hormone production, affecting approximately 2\% of the population, with higher prevalence in women, the elderly, and iodine-deficient regions. The condition results from various causes, with the most common in iodine-sufficient areas being Hashimoto's thyroiditis (chronic autoimmune thyroiditis, characterized by lymphocytic infiltration and destruction of the thyroid gland with positive anti-TPO and anti-thyroglobulin antibodies), while iodine deficiency is the most common cause worldwide; other causes include post-thyroidectomy or post-radioiodine therapy, external beam radiation, medications (lithium, amiodarone, tyrosine kinase inhibitors), central (pituitary/hypothalamic) hypothyroidism, and congenital hypothyroidism. Clinical features include fatigue, weight gain, cold intolerance, constipation, dry skin, hair loss, bradycardia, delayed relaxation of deep tendon reflexes, and depression, with severe hypothyroidism (myxedema coma) being a medical emergency with hypothermia, altered consciousness, and cardiovascular collapse. Diagnosis is by elevated TSH (primary hypothyroidism) and low free T4. Management is levothyroxine (synthetic T4) replacement, dosed to normalize TSH. Prognosis is excellent with appropriate replacement therapy.

\subsection{Thyroid Nodules and Cancer}
\label{sec:Thyroid Nodules and Cancer}
Thyroid nodules are extremely common, palpable in 4--7\% of adults and incidentally discovered on imaging in up to 50\%, with the vast majority (90--95\%) being benign. Thyroid cancer accounts for approximately 58,000 new cases annually in the US, with papillary thyroid carcinoma being the most common histological type (80\%, with excellent prognosis and 20-year survival $>$95\%), followed by follicular thyroid carcinoma (10--15\%), medullary thyroid carcinoma (5\%, arising from parafollicular C cells and producing calcitonin, associated with MEN2), and anaplastic thyroid carcinoma ($<$1\%, highly aggressive and usually fatal). The condition results from neoplastic transformation of thyroid follicular or parafollicular cells. Clinical features include a palpable thyroid nodule, sometimes with compressive symptoms. The American Thyroid Association (ATA) guidelines guide management based on ultrasound characteristics (TI-RADS classification), size, and risk factors, with fine-needle aspiration (FNA) biopsy being the standard for evaluating suspicious nodules (Bethesda classification for cytology). Management of differentiated thyroid cancer includes total thyroidectomy (or lobectomy for low-risk papillary cancer $<$4 cm), radioactive iodine ablation for intermediate/high-risk disease, and TSH suppression with levothyroxine. Prognosis is excellent for most differentiated thyroid cancers.

\subsection{Thyroiditis}
\label{sec:Thyroiditis}
Thyroiditis encompasses inflammatory conditions of the thyroid gland, with Hashimoto's thyroiditis being the most common cause of hypothyroidism in iodine-sufficient regions. Subacute (de Quervain's) thyroiditis is a painful, granulomatous thyroiditis thought to be viral in etiology, presenting with thyroid pain, tenderness, fever, and thyrotoxicosis followed by hypothyroidism and eventual recovery. Silent (painless) thyroiditis and postpartum thyroiditis are autoimmune conditions presenting with transient thyrotoxicosis followed by hypothyroidism. Acute suppurative thyroiditis is a rare bacterial infection. Drug-induced thyroiditis may be caused by amiodarone, lithium, interferon-alpha, and immune checkpoint inhibitors. The condition results from inflammation of the thyroid gland. Clinical features vary by type. Diagnosis is based on clinical presentation, thyroid function tests, and imaging. Management varies by type: NSAIDs or corticosteroids for subacute thyroiditis, levothyroxine for hypothyroid phases, beta-blockers for symptomatic thyrotoxicosis, and observation for most self-limited forms. Prognosis is generally excellent.

